10.3 The Differential of a Mapping 73
b) Show that the correspondence
$$
\Omega = \begin{pmatrix} 0 & -\omega^3 & \omega^2 \\ \omega^3 & 0 & -\omega^1 \\ -\omega^2 & \omega^1 & 0 \end{pmatrix} \leftrightarrow (\omega^1, \omega^2, \omega^3) = \omega
$$
is an isomorphism of the Lie algebras $\mathfrak{L}A_2$ and $\mathfrak{L}A_1.$
c) Verify that if the skew-symmetric matrix $\Omega$ and the vector $\omega$ correspond to
each other as shown in b), then the equality $\dot{\mathbf{r}} = [\omega, \mathbf{r}]$ holds for any vector $\mathbf{r} \in E^3,$
and the relation $P\Omega P^{-1} \leftrightarrow P\omega$ holds for any matrix $P \in SO(3).$
d) Verify that if $\mathbb{R} 
i t \rightarrow O(t) \in SO(3)$ is a smooth mapping, then the matrix
$\Omega(t) = O^{-1}(t)\dot{O}(t)$ is skew-symmetric.
e) Show that if $\mathbf{r}(t)$ is the radius vector of a point of a rotating top and $\Omega(t)$ is
the matrix $(O^{-1}(t)\dot{O}(t))$ found in d), then $\dot{\mathbf{r}}(t) = [\omega(t),\mathbf{r}(t)].$
f) Let $\mathbf{r}$ and $\omega$ be two vectors attached at the origin of $E^3.$ Suppose a right-
handed frame has been chosen in $E^3,$ and that the space undergoes a right-handed
rotation with angular velocity $|\omega|$ about the axis defined by $\omega.$ Show that $\dot{\mathbf{r}}(t) =$
$[\omega, \mathbf{r}(t)]$ in this case.
g) Summarize the results of d), e), and f) and exhibit the instantaneous angular
velocity of the rotating top discussed in Example 8.
h) Using the result of c), verify that the velocity vector $\omega$ is independent of the
choice of the fixed orthoframe in $E^3,$ that is, it is independent of the coordinate
system.
6. Let $\mathbf{r} = \mathbf{r}(s) = (x^1(s), x^2(s), x^3(s))$ be the parametric equations of a smooth
curve in $E^3,$ the parameter being arc length along the curve (the natural parametrization of the curve).
a) Show that the vector $\mathbf{e}_1(s) = \frac{d\mathbf{r}}{ds}$ tangent to the curve has unit length.
b) The vector $\frac{d\mathbf{e}_1}{ds}(s)$ is orthogonal to $\mathbf{e}_1.$ Let $\mathbf{e}_2(s)$ be the unit vector
formed from $\frac{d\mathbf{e}_1}{ds}.$ The coefficient $k(s)$ in the equality $\frac{d\mathbf{e}_1}{ds}(s) = k(s)\mathbf{e}_2(s)$ is called
the curvature of the curve at the corresponding point.
c) By constructing the vector $\mathbf{e}_3 (s) = [\mathbf{e}_1 (s), \mathbf{e}_2(s)]$ we obtain a frame $\{\mathbf{e}_1, \mathbf{e}_2, \mathbf{e}_3\}$
at each point, called the Frenet frame$^3$ or companion trihedral of the curve. Verify
the following Frenet formulas:
$$
\frac{d\mathbf{e}_1}{ds}(s) = k(s)\mathbf{e}_2(s),
$$
$$
\frac{d\mathbf{e}_2}{ds}(s) = -k(s)\mathbf{e}_1(s) + \kappa(s)\mathbf{e}_3(s),
$$
$$
\frac{d\mathbf{e}_3}{ds}(s) = -\kappa(s)\mathbf{e}_2(s).
$$
$^3$J.F. Frenet (1816–1900) – French mathematician.